%% kapa: exam=1 round=35 year=2024 period=1 form=A subject=civil_law subject_ko=민법 seq=4 number=4
%% source: Q-Net 감정평가사 기출문제 (cst00309) artlSeq=5213229 file=raw/1008/5213229/1교시 문제지 원본.pdf
%% conversion: kordoc 4.24.51 -> tools/gen1.py (auto); review=proofread-sonnet; flags=none
\begin{question}{4}
제한능력자에 관한 설명으로 옳은 것은?
\begin{choices}
  \item 미성년자가 법정대리인으로부터 허락을 얻은 특정한 영업에 관해서는 법정대리인의 대리권이 소멸한다.
  \item 제한능력을 이유로 하는 취소는 특별한 사정이 없는 한 선의의 제3자에게 대항할 수 없다.
  \item 제한능력자의 단독행위는 유효한 추인이 있은 후에도 상대방이 거절할 수 있다.
  \item 가정법원은 취소할 수 없는 피성년후견인의 법률행위의 범위를 정할 수 없다.
  \item 가정법원은 정신적 제약으로 특정한 사무에 관해 후원이 필요한 사람에 대해서는 본인의 의사에 반하더라도 특정후견 심판을 할 수 있다.
\end{choices}
\end{question}
\answer{1}
